\begin{remark}[7.3.11]
\label{IV.2.s7.3.11-fr}
Dejamos al lector la tarea de formular la propiedad $\boldsymbol R$
correspondiente a cada uno de los ejemplos de (7.3.8) y de verificar,
en cada caso, las condiciones $(\mathrm R_{\mathrm I})$ y
$(\mathrm R_{\mathrm{II}})$ con ayuda de los resultados del \S\,6.
De hecho, salvo en el ejemplo (i) de (7.3.8), la propiedad
$\boldsymbol R$ satisface en los demás casos la condición siguiente:
\begin{enumerate}
\item[$(\mathrm R'_{\mathrm I})$] Sean $X$ e $Y$ dos preesquemas
localmente noetherianos y $f:X\to Y$ un $\boldsymbol P$-\emph{morfismo}
(para la propiedad $\boldsymbol P$ del primer o del segundo tipo
definida por $\boldsymbol R$); entonces, para todo $x\in X$, la
propiedad $\boldsymbol R(\mathscr O_{f(x)})$ implica la propiedad
$\boldsymbol R(\mathscr O_x)$.
\end{enumerate}
Los razonamientos de (6.6.1) y (6.8.3) demuestran que, si
$\boldsymbol R$ satisface las condiciones $(\mathrm R'_{\mathrm I})$
y $(\mathrm R_{\mathrm{II}})$, entonces $\boldsymbol P$ satisface las
condiciones $(\mathrm P'_{\mathrm I})$ (7.3.5, (iii)) y
$(\mathrm P_{\mathrm{II}})$.
\end{remark}
